\documentclass[10pt, aspectratio=169]{beamer}
\usepackage{../beamer_theme_408}

\title{408考研零基础 --- 计算机组成原理}
\subtitle{2-13 CPU概述}
\author{}
\date{}

\begin{document}

\frame{\titlepage}

\begin{frame}{本节目标}
    \begin{enumerate}
        \item 掌握CPU的四大基本功能
        \item 区分通用寄存器与专用寄存器
        \item 理解运算器的组成与作用
        \item 了解控制器的基本工作流程
    \end{enumerate}
\end{frame}

\begin{frame}{CPU的基本功能}
    \begin{center}
        \textbf{CPU = 控制器 + 运算器}
    \end{center}
    \vspace{0.5em}
    \begin{enumerate}
        \item \textbf{指令控制} —— 取指令、分析指令、执行指令
        \item \textbf{操作控制} —— 产生控制信号，控制各部件协调工作
        \item \textbf{时间控制} —— 为每条指令提供时序控制
        \item \textbf{数据加工} —— 对数据进行算术运算和逻辑运算
    \end{enumerate}
    \vspace{0.5em}
    \begin{block}{延伸}
        此外CPU还负责中断处理、异常处理等控制功能。
    \end{block}
\end{frame}

\begin{frame}{CPU的内部结构}
    \begin{center}
        \begin{tabular}{c}
            \hline
            \textbf{CPU} \\
            \hline
            \begin{tabular}{c}
                \textbf{控制器} \\
                \begin{tabular}{c}
                    \textbullet\ 取指令 \\
                    \textbullet\ 译码 \\
                    \textbullet\ 发控制信号
                \end{tabular}
            \end{tabular}
            \quad
            \begin{tabular}{c}
                \textbf{运算器} \\
                \begin{tabular}{c}
                    \textbullet\ 算术运算 \\
                    \textbullet\ 逻辑运算 \\
                    \textbullet\ 移位
                \end{tabular}
            \end{tabular}
            \\
            \hline
            \multicolumn{1}{c}{\textbf{寄存器组}} \\
            \hline
        \end{tabular}
    \end{center}
    \vspace{0.3em}
    \begin{itemize}
        \item 控制器是"指挥中心"，运算器是"加工厂"
        \item 寄存器组是中间数据暂存区
    \end{itemize}
\end{frame}

\begin{frame}{寄存器概述}
    \textbf{寄存器}：CPU内部的高速存储单元，读写速度远高于主存
    \vspace{0.3em}
    \begin{center}
        \begin{tabular}{|c|c|c|}
            \hline
            \textbf{分类} & \textbf{名称} & \textbf{功能} \\
            \hline
            通用 & 通用寄存器 & 存放操作数、中间结果 \\
            专用 & PC & 程序计数器，存下一条指令地址 \\
            专用 & IR & 指令寄存器，存当前指令 \\
            专用 & MAR & 存储器地址寄存器 \\
            专用 & MDR & 存储器数据寄存器 \\
            专用 & PSW & 程序状态字 / 标志寄存器 \\
            \hline
        \end{tabular}
    \end{center}
\end{frame}

\begin{frame}{专用寄存器详解}
    \begin{itemize}
        \item \textbf{PC（程序计数器）}
            \begin{itemize}
                \item 存放下一条指令在主存中的地址
                \item 取指后自动加1（或加指令字节数）
                \item 转移指令会修改PC值
            \end{itemize}
        \item \textbf{IR（指令寄存器）}
            \begin{itemize}
                \item 存放当前正在执行的指令
                \item 控制器根据IR内容译码产生控制信号
            \end{itemize}
        \item \textbf{MAR（存储器地址寄存器）}
            \begin{itemize}
                \item 存放要访问的主存单元的地址
                \item 地址线宽度决定最大寻址空间
            \end{itemize}
        \item \textbf{MDR（存储器数据寄存器）}
            \begin{itemize}
                \item 存放从主存读出或写入主存的数据
            \end{itemize}
        \item \textbf{PSW（程序状态字）}
            \begin{itemize}
                \item 存放运算结果标志（ZF、OF、SF、CF等）
                \item 存放中断标志、CPU运行模式等
            \end{itemize}
    \end{itemize}
\end{frame}

\begin{frame}{运算器的组成}
    \textbf{运算器（ALU）}的核心部件：
    \vspace{0.3em}
    \begin{enumerate}
        \item \textbf{ALU（算术逻辑单元）}
            \begin{itemize}
                \item 执行加、减、与、或等基本运算
                \item 组合逻辑电路
            \end{itemize}
        \item \textbf{寄存器组}
            \begin{itemize}
                \item 暂存操作数和运算结果
                \item 加速数据存取
            \end{itemize}
        \item \textbf{数据通路}
            \begin{itemize}
                \item 连接各寄存器和ALU的线路
                \item 由多路选择器、三态门控制
            \end{itemize}
        \item \textbf{标志寄存器（PSW的一部分）}
            \begin{itemize}
                \item 记录运算结果的特性
            \end{itemize}
    \end{enumerate}
\end{frame}

\begin{frame}{控制器的功能与组成}
    \textbf{控制器功能}：
    \begin{enumerate}
        \item \textbf{取指令}：从主存中取出指令送IR
        \item \textbf{分析指令}：译码器分析IR的操作码
        \item \textbf{执行指令}：发控制信号驱动各部件
        \item \textbf{处理异常/中断}
    \end{enumerate}
    \vspace{0.3em}
    \textbf{控制器组成}：
    \begin{itemize}
        \item 程序计数器 PC
        \item 指令寄存器 IR
        \item 指令译码器
        \item 时序控制逻辑
        \item 微操作信号发生器
    \end{itemize}
\end{frame}

\begin{frame}{CPU的工作流程}
    \begin{center}
        \textbf{取指令 $\rightarrow$ 译码 $\rightarrow$ 执行 $\rightarrow$ 中断检查}
    \end{center}
    \vspace{0.5em}
    \begin{enumerate}
        \item \textbf{取指令}
            \begin{itemize}
                \item PC $\rightarrow$ MAR $\rightarrow$ 主存 $\rightarrow$ MDR $\rightarrow$ IR
                \item PC $\leftarrow$ PC + 1
            \end{itemize}
        \item \textbf{译码}
            \begin{itemize}
                \item 控制器分析IR中的操作码
                \item 生成相应的微操作控制信号
            \end{itemize}
        \item \textbf{执行}
            \begin{itemize}
                \item 根据控制信号进行运算或数据传送
                \item 结果写回目标寄存器或内存
            \end{itemize}
        \item \textbf{检查中断}
            \begin{itemize}
                \item 若有中断请求，进入中断处理
                \item 否则进入下一条指令的取指
            \end{itemize}
    \end{enumerate}
\end{frame}

\begin{frame}{CPU的主要性能指标}
    \begin{itemize}
        \item \textbf{主频}：CPU内部时钟频率，单位GHz（例：3.0 GHz）
        \item \textbf{CPI（Clock cycles Per Instruction）}
            \begin{itemize}
                \item 执行一条指令所需的时钟周期数
                \item 不同指令CPI不同
            \end{itemize}
        \item \textbf{IPC（Instructions Per Cycle）}
            \begin{itemize}
                \item 每个时钟周期执行的指令数 = $1/\text{CPI}$
            \end{itemize}
        \item \textbf{MIPS（Million Instructions Per Second）}
            \begin{itemize}
                \item $\text{MIPS} = \dfrac{\text{主频}}{\text{CPI} \times 10^6}$
            \end{itemize}
    \end{itemize}
    \vspace{0.3em}
    \begin{block}{关系}
        主频高不一定性能好，还需考虑CPI和指令集效率。
    \end{block}
\end{frame}

\begin{frame}{本节总结}
    \begin{enumerate}
        \item CPU功能：指令控制、操作控制、时间控制、数据加工
        \item 寄存器分通用和专用，PC/IR/MAR/MDR/PSW是核心专用寄存器
        \item 运算器核心是ALU，配合寄存器组和数据通路完成运算
        \item 控制器负责取指、译码、发控制信号，协调指令执行
    \end{enumerate}
    \vspace{1em}
    \begin{center}
        \textcolor{blue}{\textbf{下节预告：2-14 指令执行过程 --- 一条指令的旅程}}
    \end{center}
\end{frame}

\end{document}
